\chapter{Теоретические основы исследования}
\section{Понятие и сущность эмпатии}
Эмпатия -- это ощущение сходства между чувствами субъекта и чувствами других людей. Ученые в области когнитивных нейронаук отмечают, что исследование подобных эмпатии субъективных психологических феноменов возможно только с помощью мультидисциплинарного подхода \cite{Decety2004}. 

Данный подход позволяет лучше понять механизмы обработки информации, которые порождают это субъективное психологическое явление \cite{Uysal2019}. Мультидисциплинарные исследования включают в себя когнитивную нейронауку, социальную психологию, клиническую нейропсихологию и другие области, что позволяет более полно исследовать различные аспекты эмпатии \cite{Duan2018}.

Выделяют два основных аспекта эмпатии: когнитивный и аффективный аспекты. Функции когнитивной эмпатии отвечают за осознание чувств другого человека вместе с пониманием причин этих чувств \cite{Moscrop2001}. Функции аффективной эмпатии отвечают за агентное поведение -- способность эмоционально и активно откликаться и сопереживать \cite{Kaluzeviciute2020}.

Эмпатия играет значительную роль в принятии решений, особенно в контексте моральных и этических дилемм. Исследования показывают, что эмпатия может влиять на эмоциональные реакции и восприятие моральных ситуаций, что, в свою очередь, может модифицировать процесс принятия решений. Например, более эмпатичные люди склонны воспринимать моральные дилеммы как более неприятные и вызывающие большее эмоциональное возбуждение, что может повлиять на их решения~\cite{Cecchetto2017}. 

К. Даниэль Батсон предложил теорию эмпатии, которая выделяет два качественно различных вида сопереживания: личное беспокойство и эмпатию~\cite{Batson1987}. Эти два вида эмоций имеют разные мотивационные последствия.

Личное беспокойство рассматривается как чувство тревоги, беспокойства и расстройства, возникающее при восприятии потребности другого человека. Личное беспокойство вызывает эгоистическую мотивацию, направленную на уменьшение собственного неприятного состояния.

Эмпатия рассматривается как чувство симпатии, сострадания и нежности, возникающее при восприятии потребностей другого человека. Эмпатия вызывает альтруистическую мотивацию, направленную на снижение неудовлетворенности другого человека, а не на решение собственных проблем.

Следует отметить, что на эмпатию разное влияние оказывают как положительные, так и отрицательные эмоции. 

Положительные эмоции могут усиливать субъективное восприятие эмпатии, но не всегда приводят к улучшению восприятия эмпатии. Например, люди, испытывающие положительные  эмоций могут лучше распознавать положительные эмоции у других, но хуже справляются с распознаванием интенсивных негативных эмоций \cite{Devlin2014}.

Отрицательные эмоции могут снижать способность к распознаванию эмоций и эмпатии. Например, личный стресс, вызванный негативными эмоциями, отрицательно влияет на точность распознавания эмпатии~\cite{Israelashvili2020}.

Положительная эмпатия связана с повседневными просоциальными действиями, направленными на улучшение настроения других, в то же время отрицательная эмпатия в большей степени связана с оказанием помощи в ситуациях, когда другие люди испытывают острую необходимость в поддержке~\cite{Andreychik2015}.


Способность понимать и предсказывать ментальные состояния других людей, такие как убеждения, желания и намерения, позволяет Теория разума (англ. Theory of Mind, ToM). Исследования показывают, что ToM играет ключевую роль в развитии эмпатии и социального понимания~\cite{Brown2017}. Теория разума позволяет людям не только понимать эмоции других, но и предсказывать их поведение, что важно для стратегического достижения поставленных целей~\cite{Artinger2014}. 

Исследования показывают, что теория разума (ToM) и эмпатия играют важную роль в выравнивании целей ИИ с человеческими ценностями. Эмпатия, как утверждается, необходима для этического поведения ИИ, несмотря на то, что часто игнорируется в пользу более дедуктивных подходов \cite{https://doi.org/10.48550/arxiv.2312.07579}.


\section{Учет эмпатии в больших языковых моделях}

Большие языковые модели (БЯМ) представляют собой современные системы искусственного интеллекта, предназначенные для понимания и генерации естественного языка посредством обучения на обширных массивах текстовых данных. Эти модели, к числу которых относятся GPT, LLaMA и Gemini, включают в себя миллиарды параметров, что обеспечивает им возможность выполнять широкий спектр задач по обработке естественного языка с высокой точностью и эффективностью~\cite{https://doi.org/10.48550/arxiv.2402.06196}.

Для понимания механизмов ответа БЯМ с учетом эмпатии необходимо механизм внутренней работы этих моделей.  К методам понимания внутренней работы больших языковых моделей относятся механическая интепретируемость и обучение с подкреплением на основе человеческих предпочтений.

Механистическая интерпретируемость в машинном обучении направлена на понимание внутренних механизмов моделей путем анализа их компонентов и взаимодействий. Этот подход стремится объяснить поведение моделей через их внутренние структуры, такие как веса и активации нейронных сетей, и выявить, как эти структуры способствуют решению задач \cite{https://doi.org/10.48550/arxiv.2309.16042}.


Механическая интерпретируемость важна для обеспечения прозрачности и доверия к системам искусственного интеллекта -- она позволяет понять, как внутренние механизмы моделей влияют на их способность распознавать и генерировать эмоции. В исследовании \cite{
https://doi.org/10.48550/arxiv.2403.03121} было показано, что большие языковые модели могут перенимать гендерные эмоциональные стереотипы от людей, что связано с их внутренними механизмами и обучением на данных, содержащих такие стереотипы.

Кроме того, интеграция человеческих эмоций и этических соображений в LLM требует глубокого понимания их внутренних механизмов. Это позволяет моделям не только генерировать необходимые эмоции, но и следовать этическим нормам~\cite{https://doi.org/10.48550/arxiv.2404.13071}.

Обучение с подкреплением на основе обратной связи от человека (англ. Reinforcement learning from human feedback) используется для настройки LLM в соответствии с человеческими предпочтениями и ценностями, что включает в cебя в том числе и эмпатию. RLHF позволяет моделям генерировать ответы, которые более соответствуют ожиданиям и эмоциональным реакциям людей.

Также RLHF позволяет моделям контролировать эмоциональный тон своих ответов. Например, метод Direct Preference Optimization (DPO) улучшает способность моделей контролировать эмоциональную глубину ответа~\cite{https://doi.org/10.48550/arxiv.2305.18290}.

Эмпатию в больших языковых моделях (LLM) оценивают с помощью различных методов, включая субъективные оценки людьми и объективные метрики.

Чаще всего эмпатию оценивают с помощью человеческих оценок. Например, в одном исследовании люди оценивали ответы моделей на основе их восприятия эмпатии в ответах на жизненные ситуации, такие как рабочие проблемы или семейные отношения. Результаты показали, что ответы LLM часто воспринимаются как более эмпатичные, чем ответы людей~\cite{lee2024large}.

В редких случаях используются объективные метрики для оценки эмпатии. Например, Пател и Фан использовали шкалы, такие как Toronto Alexithymia Scale (TAS-20) и Empathy Quotient (EQ-60), для оценки способности моделей распознавать и описывать эмоции.  В их исследовании были оценены уровни эмпатии и способность к идентификации и описанию эмоций у трех языковых моделей: Bard, GPT-3.5 и GPT-4. Результаты показали, что менее совершенные модели (Bard и GPT-3.5) показали худшие результаты на TAS-20, что близко к алекситимии, состоянию, характеризующемуся значительными трудностями в распознавании, выражении и описании эмоций. Однако GPT-4 достигла уровня, близкого к человеческому, с двумя подкатегориями, превосходящими человеческие показатели~\cite{Patel2023}.

Оценивая эмпатию, необходимо учитывать агентное поведение. Агентное поведение относится к действиям и реакциям, которые искусственные агенты (например, роботы или виртуальные ассистенты) демонстрируют в ответ на взаимодействие с пользователями.

Современные алгоритмы машинного обучения позволяют агентам анализировать и интерпретировать мультимодальные сигналы, такие как голос, мимика и другие сенсорные входы, для создания адекватных эмпатических ответов. Например, глубокое обучение и нейронные сети могут быть использованы для распознавания эмоциональных состояний и адаптации поведения агентов в реальном времени.

Агентное поведение в контексте эмпатии также является важным аспектом, особенно в исследованиях, связанных с искусственными агентами и роботами. Например, исследования показывают, что эмпатическое поведение агентов может улучшить взаимодействие с пользователями и повысить их автономность~\cite{10.11453570945.3607291}. Это включает в себя использование мультимодальных сигналов, таких как голос, мимика и другие сенсорные входы, для создания адекватных эмпатических ответов~\cite{10.1145/2912150}.

Агентное поведение в виртуальных средах оказывает значительное влияние на восприятие эмпатии. Исследования показывают, что различные невербальные маркеры, такие как направление взгляда, выражение лица и наклон туловища, могут существенно влиять на восприятие эмпатии от виртуальных агентов~\cite{doi:10.1080/10447318.2021.1987680}. Например, выражение боли на лице агента воспринимается как более эмпатичное, чем нейтральное выражение, независимо от других факторов.  Агенты, использующие мультимодальные входные данные для интерпретации поведения пользователя и генерации эмпатических ответов, воспринимаются как более эмпатичные и улучшают вовлеченность пользователей.


\section{Большие языковые модели в контексте теории игр}

Теория игр -- это математическое исследование взаимодействий между независимыми, рациональными и самозаинтересованными агентами \cite{LeytonBrown2008}. Она используется для анализа стратегических взаимодействий среди рациональных лиц, принимающих решения, и охватывает такие ключевые компоненты, как игроки, стратегии и выигрыши.

Основные концепции теории игр включают в себя нормальные и расширенные формы игр, а также классические примеры, такие как дилемма заключенного и равновесие Нэша. Теория игр применяется в различных областях, включая экономику, политологию, биологию и психологию, что делает её универсальным инструментом для понимания человеческого поведения и принятия решений.
 
Концепция равновесия Нэша позволяет предсказывать устойчивые состояния, в которых ни один из участников не имеет стимула изменять свою стратегию. Это особенно полезно для анализа рынков с несовершенной конкуренцией~\cite{Kumar2024}.

Виды игр в теории игр представлены в таблице~\ref{Namelists}.

\begin{longtable}[c]{|p{0.3\textwidth}|p{0.65\textwidth}|} %longtable* появляется из пакета ltcaption и даёт ненумерованную таблицу
    \caption{Виды игр}\label{Namelists}
    \\
    \hline
    %\multicolumn{4}{|c|}{\textbf{Файл puma\_namelist}}        \\ \hline
    Игра & Описание               \\ \hline
    \endfirsthead   
    \multicolumn{2}{l}{\hspace{-2mm}Продолжение таблицы~\ref{Namelists}}        \\ %\hline
\hline
    Игра & Описание              \\ \hline \endhead        \hline % \multicolumn{4}{|c|}{\small\slshape (окончание)}        \\ \hline
    % Параметр & Умолч. & Тип & Описание               \\ \hline
    %                                             \endlasthead        \hline
    %%\multicolumn{2}{|r|}{\small\slshape продолжение следует}  \\ \hline
    \endfoot        \hline
    \endlastfoot
	Статистические игры с полной информацией & В этих играх все игроки принимают решения одновременно, и каждый игрок знает стратегии и выигрыши всех остальных игроков. Примером является равновесие Нэша. \\
     \hline
	Динамические игры с полной информацией & Эти игры включают последовательные ходы, где игроки принимают решения по очереди. Примером является субигровое совершенное равновесие. \\
     \hline
	Статические игры с неполной информацией & В этих играх игроки не знают стратегий и выигрышей других игроков. Примером является байесовское равновесие Нэша. \\
     \hline
Динамические игры с неполной информацией & Эти игры включают последовательные ходы и неполную информацию о стратегиях и выигрышах других игроков. Примером является совершенное байесовское равновесие. \\
\hline

Дифференциальные игры & Эти игры используются для моделирования динамических процессов в экономике, таких как накопление капитала и конкуренция на олигополистических рынках. \\
\hline
\end{longtable}


В зависимости от выбранной стратегии или от распределения выигрыша можно выделить другую классификацию игр:

\begin{enumerate}[label=\arabic*)] % https://tex.stackexchange.com/a/476052/104425
\item В некооперативных играх основное понятие -- равновесие Нэша, где каждый игрок выбирает стратегию, максимизирующую его выигрыш, учитывая стратегии других игроков.

\item В кооперативных играх основное понятие -- ядро, которое представляет собой множество распределений выигрышей, которые не могут быть улучшены никакой коалицией.

\item В играх с неполной информацией основное понятие -- байесовское равновесие, где игроки формируют ожидания о стратегиях других игроков на основе вероятностных распределений.

\item В симметричных играх основное понятие -- симметричное равновесие, где все игроки выбирают одинаковые стратегии.
\end{enumerate}

Основные характеристики типов игр в теории игр представлены в таблице~\ref{gamechar}. 

\begin{longtable}[c]{|p{0.21\textwidth}|p{0.29\textwidth}|p{0.42\textwidth}|}
    \caption{Характеристика игр}\label{gamechar}
    \\
    \hline
    %\multicolumn{4}{|c|}{\textbf{Файл puma\_namelist}}        \\ \hline
    Тип игры & Определение & Пример             \\ \hline

    \endfirsthead  
    \multicolumn{3}{l}{\hspace{-2mm}Продолжение таблицы~\ref{gamechar}}        \\ %\hline
 \hline
    Тип игры & Определение & Пример             \\ \hline
    \hline \endhead  
    % Параметр & Умолч. & Тип & Описание               \\ \hline
    %                                             \endlasthead        \hline

 %   \multicolumn{3}{|r|}{\small\slshape продолжение следует}  \\ \hline
    \endfoot        \hline
    \endlastfoot
Неко\-о\-пе\-ра\-тив\-ные & Игры, в которых каждый участник действует независимо, без кооперации или коммуникации с другими игроками. & Дилемма заключенного, где два игрока должны решить, сотрудничать или предать друг друга, не зная выбора другого. Игра <<Диктатор>>, где диктатор имеет полную власть над распределением ресурсов, в то время как другой игрок (получатель) не может повлиять на это решение.  \\
\hline
Ко\-о\-пе\-ра\-тив\-ны\-е  & Игры, в которых игроки могут формировать коалиции и совместно принимать решения для достижения общего выигрыша. & Игры по дележу ресурсов, где игроки формируют коалиции для максимизации общего выигрыша (игра <<Общественное благо>>). \\ \hline

Э\-во\-лю\-ци\-он\-ные & Игры, моделирующие динамику стратегий в популяциях, где стратегии игроков могут изменяться со временем под воздействием естественного отбора. &  Игра <<Ястреб и Голубь>>, моделирующая поведение агрессивных и мирных стратегий в популяции. \\ \hline

С непол\-ной ин\-фор\-мац\-и\-ей & Игры, в которых игроки не обладают полной информацией о состоянии игры или стратегиях других игроков. & Байесовские игры, где игроки имеют неполную информацию о стратегиях и предпочтениях других игроков. \\ \hline

Стa\-ти\-чес\-кие и ди\-на\-ми\-чес\-кие & Статические игры рассматривают однократные взаимодействия, тогда как динамические игры включают многократные взаимодействия с возможностью изменения стратегий & Статическая игра <<Камень, ножницы, бумага>> и динамическая игра <<Повторяющаяся дилемма заключенного>>.\\ \hline

Симметричные & Игры, в которых стратегии и выигрыши игроков симметричны. & Игра <<Охота на оленя>>, где оба игрока выигрывают, если сотрудничают, и проигрывают, если один из них предает.  \\ \hline

\end{longtable}

Существуют определенные ограничения применений LLM в теории игр. К таким ограничениям относятся ограниченные способности к стратегическому мышлению: LLM демонстрируют различные поведения в зависимости от игровых сценариев. Например, они плохо справляются с играми с полной и детерминированной информацией, но могут быть конкурентоспособными в вероятностных играх.  Даже современные LLM, такие как GPT-4, показывают значительные различия по сравнению с людьми в аспектах рациональности. Они могут не строить желания на основе необычных предпочтений и не всегда корректно интерпретируют простые паттерны.  Другие модели, такие как LLaMa-2, лучше адаптируются к структуре игры и контексту, но все равно имеют ограничения~\cite{duan2024gtbench}.

Для преодоления этих трудностей предлагаются различные подходы. Например, метод <<K-Level Reasoning>> улучшает предсказание последующих ходов соперников и стратегическое принятие решений~\cite{zhang2024klevel}. Также использование внутренней мотивации, как в методе Motif, позволяет агентам достигать высоких результатов, даже без прямого взаимодействия с окружающей средой~\cite{klissarov2023motif}.

В исследовании, проведенным Квоном, использовались текстовые подсказки для создания прокси-функций вознаграждения в играх, таких как «Ультиматум». Модели GPT-3 использовались для оценки поведения агентов и выдачи соответствующих сигналов вознаграждения, что позволило обучить агентов, соответствующих целям пользователя~\cite{kwon2023reward}. 

С.~Фэлпсом изучались и способности GPT-3.5 к кооперативному, конкурентному, альтруистическому и эгоистичному поведению в социальных дилеммах, включая игры «Ультиматум» и «Диктатор». Результаты показали, что модели могут переводить описания поведения в соответствующие действия, но испытывают трудности с адаптацией на основе условной взаимности~\cite{phelps2023investigating}. Результаты подтверждаются исследованием Э.~Акаты: в исследовании изучалось поведение моделей GPT-3, GPT-3.5 и GPT-4 в повторяющихся играх. Результаты показали, что модели хорошо справляются с задачами, где выгодно учитывать собственные интересы, но испытывают трудности в играх, требующих координации~\cite{akata2023playing}, при этом модели могут воспроизводить распределения ответов, согласующиеся с предыдущими экспериментальными данными.

LLM часто демонстрируют предпочтение к определенным стратегиям эмоциональной поддержки, что снижает их эффективность. Внешняя помощь может уменьшить эту предвзятость, но LLM все еще не могут стать хорошими эмоциональными поддерживающими агентами без дополнительной помощи. Так, иследование Д.~Дзяо показало, что культурные ценности индивидуализма и коллективизма влияют на поведение в играх «Ультиматум» и «Диктатор». Участники, подвергшиеся воздействию коллективистских текстов, демонстрировали более альтруистичное поведение~\cite{Jiao2023}.


Существующие бенчмарки для оценки эмоционального интеллекта LLM имеют ограничения, такие как фокус на распознавании эмоций и наличие ошибок в аннотациях. EmoBench предлагает более комплексный подход, включающий понимание и применение эмоций. Исследование C.~Сабура оценивает эмоциональный интеллект LLM, включая распознавание, интерпретацию и понимание эмоций. Некоторые модели достигают высоких результатов, но механизмы достижения этих результатов отличаются от человеческих~\cite{sabour2024}.

LLM сталкиваются с трудностями и в детекции тонких эмоций в тексте из-за нескольких факторов. Во-первых, такие модели часто переобучаются на случайные корреляции, такие как позиционная предвзятость, что снижает их способность точно классифицировать эмоции. Во-вторых, LLMs могут испытывать трудности с интерпретацией контекста и нюансов, что особенно важно для распознавания тонких эмоций. Например, LLM отражают гендерные стереотипы в атрибуции эмоций, что приводит к предвзятости в генерации эмоциональных ответов. Например, женщины чаще ассоциируются с эмпатией, а мужчины — с гневом~\cite{https://doi.org/10.48550/arxiv.2403.03121}.


\FloatBarrier
